% questions.tex —— 2014 年·四川卷·理科 · 题目内容片段（无答案）
% 由 tools/generate_projects.py 依来源总账自 2014/2014_数学真题_四川理_试题.md 生成，勿手改（改源 MD 后重跑）


\begin{gknotice}
\textbf{来源说明}：本题面依据固定源 PDF 的两页 A3 题面逐题转录，原卷印刷页码为 919—920。第 919 面为选择题 1—10、填空题 11—15 与解答题 16—17；第 920 面为解答题 18—21。全卷三处节头依次为「一、选择题」「二、填空题」「三、解答题」；题号与索引一致。

\textbf{解析归属}：固定源不附答案或解析；配套答案与解析均由本库按本题面独立复算。

\textbf{图形说明}：第 5 题的程序框图、第 8 题的正方体、第 13 题的气球测河宽示意图、第 18 题的三棱锥及侧视图与俯视图由项目内 TikZ 矢量源重绘；其余各题原页无图，第 20 题椭圆亦无图，均不补绘。
\end{gknotice}

\gkpart{一、选择题}

\begin{question}
\gkti{1}已知集合 $A=\{x\mid x^2-x-2\leqslant 0\}$，集合 $B$ 为整数集，则 $A\cap B=$\mbox{（　　）}

（A）$\{-1,0,1,2\}$；（B）$\{-2,-1,0,1\}$；（C）$\{0,1\}$；（D）$\{-1,0\}$。

\end{question}
\begin{question}
\gkti{2}在 $x(1+x)^6$ 的展开式中，含 $x^3$ 项的系数为\mbox{（　　）}

（A）$30$；（B）$20$；（C）$15$；（D）$10$。

\end{question}
\begin{question}
\gkti{3}为了得到函数 $y=\sin(2x+1)$ 的图象，只需把函数 $y=\sin 2x$ 的图象上所有的点\mbox{（　　）}

（A）向左平行移动 $\dfrac12$ 个单位长度；（B）向右平行移动 $\dfrac12$ 个单位长度；（C）向左平行移动 $1$ 个单位长度；（D）向右平行移动 $1$ 个单位长度。

\end{question}
\begin{question}
\gkti{4}若 $a>b>0$，$c<d<0$，则一定有\mbox{（　　）}

（A）$\dfrac ac>\dfrac bd$；（B）$\dfrac ac<\dfrac bd$；（C）$\dfrac ad>\dfrac bc$；（D）$\dfrac ad<\dfrac bc$。

\Needspace{0.4\textheight}
\end{question}
\begin{question}
\gkti{5}执行如图所示的程序框图，如果输入的 $x$，$y\in\mathbf R$，则输出的 $S$ 的最大值为\mbox{（　　）}

（A）$0$；（B）$1$；（C）$2$；（D）$3$。
\nopagebreak[4]
\begin{center}% ============================================================
% f2014_sichuan_science_q05_flowchart.tex
% 2014 四川卷（理科）第 5 题：程序框图
% 来源 PDF：普通高考/2014/2014四川理.pdf
% 原页：第 1 页（A3 横排）左栏下部，印面 919，第 5 题题干右下
% 结构：开始 → 输入 x,y → 判断「x≥0，y≥0，x+y≤1？」（菱形）；
%       是：S=2x+y；否：S=1；两路汇合 → 输出 S → 结束。
% ============================================================
\begin{tikzpicture}[>=Stealth, line width=0.7pt, font=\small,
  x=1cm, y=1cm, line cap=round, line join=round]

  % 开始
  \draw[rounded corners=6pt] (-0.95,0.45) rectangle (0.95,-0.45);
  \node at (0,0) {开始};
  \draw[->] (0,-0.45) -- (0,-0.95);

  % 输入
  \draw (-1.25,-0.95) -- (1.75,-0.95) -- (1.25,-1.65) -- (-1.75,-1.65) -- cycle;
  \node at (0,-1.30) {输入 $x$，$y$};
  \draw[->] (0,-1.65) -- (0,-2.15);

  % 判断（菱形，与原页一致）
  \draw (-2.95,-3.15) -- (0,-2.15) -- (2.95,-3.15) -- (0,-4.15) -- cycle;
  \node at (0,-2.95) {$x\geqslant 0$，$y\geqslant 0$，};
  \node at (0,-3.38) {$x+y\leqslant 1$？};

  % 是：左支
  \draw[->] (-2.95,-3.15) -- (-4.15,-3.15) -- (-4.15,-4.30);
  \node at (-3.35,-2.83) {是};
  \draw (-5.45,-4.30) rectangle (-2.85,-5.10);
  \node at (-4.15,-4.70) {$S=2x+y$};

  % 否：右支
  \draw[->] (2.95,-3.15) -- (4.15,-3.15) -- (4.15,-4.30);
  \node at (3.35,-2.83) {否};
  \draw (2.85,-4.30) rectangle (5.45,-5.10);
  \node at (4.15,-4.70) {$S=1$};

  % 两路汇合
  \draw (-4.15,-5.10) -- (-4.15,-5.50) -- (4.15,-5.50) -- (4.15,-5.10);
  \draw[->] (0,-5.50) -- (0,-5.85);

  % 输出
  \draw (-1.15,-5.85) -- (1.85,-5.85) -- (1.35,-6.55) -- (-1.65,-6.55) -- cycle;
  \node at (0.10,-6.20) {输出 $S$};
  \draw[->] (0.10,-6.55) -- (0.10,-6.95);

  % 结束
  \draw[rounded corners=6pt] (-0.85,-6.95) rectangle (1.05,-7.85);
  \node at (0.10,-7.40) {结束};
\end{tikzpicture}
\end{center}
\nopagebreak[4]
\end{question}
\begin{question}
\gkti{6}六个人从左至右排成一行，最左端只能排甲或乙，最右端不能排甲，则不同的排法共有\mbox{（　　）}

（A）$192$ 种；（B）$216$ 种；（C）$240$ 种；（D）$288$ 种。

\end{question}
\begin{question}
\gkti{7}平面向量 $\boldsymbol a=(1,2)$，$\boldsymbol b=(4,2)$，$\boldsymbol c=m\boldsymbol a+\boldsymbol b$（$m\in\mathbf R$），且 $\boldsymbol c$ 与 $\boldsymbol a$ 的夹角等于 $\boldsymbol c$ 与 $\boldsymbol b$ 的夹角，则 $m=$\mbox{（　　）}

（A）$-2$；（B）$-1$；（C）$1$；（D）$2$。

\Needspace{0.4\textheight}
\end{question}
\begin{question}
\gkti{8}如图，在正方体 $ABCD-A_1B_1C_1D_1$ 中，点 $O$ 为线段 $BD$ 的中点。设点 $P$ 在线段 $CC_1$ 上，直线 $OP$ 与平面 $A_1BD$ 所成的角为 $\alpha$，则 $\sin\alpha$ 的取值范围是\mbox{（　　）}

（A）$\left[\dfrac{\sqrt3}3,1\right]$；（B）$\left[\dfrac{\sqrt6}3,1\right]$；（C）$\left[\dfrac{\sqrt6}3,\dfrac{2\sqrt2}3\right]$；（D）$\left[\dfrac{2\sqrt2}3,1\right]$。
\nopagebreak[4]
\begin{center}% ============================================================
% f2014_sichuan_science_q08_cube.tex
% 2014 四川卷（理科）第 8 题：正方体中的线面角
% 来源 PDF：普通高考/2014/2014四川理.pdf
% 原页：第 1 页（A3 横排）中栏上部，印面 919，第 8 题题干上方
% 说明：正方体 ABCD-A1B1C1D1，O 为 BD 中点（实心点）；图中画出
%       截面三角形 A1BD（A1B 为实线，A1D、BD 为虚线）。
% ============================================================
\begin{tikzpicture}[>=Stealth, line width=0.7pt, font=\small,
  x=1cm, y=1cm, line cap=round, line join=round]

  \coordinate (A)  at (0,0);
  \coordinate (B)  at (2,0);
  \coordinate (C)  at (2.86,0.60);
  \coordinate (D)  at (0.86,0.60);
  \coordinate (A1) at (0,2);
  \coordinate (B1) at (2,2);
  \coordinate (C1) at (2.86,2.60);
  \coordinate (D1) at (0.86,2.60);
  \coordinate (O)  at ($(B)!0.5!(D)$);

  % 可见棱
  \draw (A) -- (B) -- (C) -- (C1) -- (D1) -- (A1) -- cycle;
  \draw (A) -- (A1);
  \draw (B) -- (B1);
  \draw (B1) -- (C1);
  \draw (A1) -- (B1);

  % 不可见棱
  \draw[dashed] (A) -- (D);
  \draw[dashed] (D) -- (C);
  \draw[dashed] (D) -- (D1);

  % 截面三角形 A1BD
  \draw (A1) -- (B);
  \draw[dashed] (A1) -- (D);
  \draw[dashed] (B) -- (D);

  % 点 O
  \fill (O) circle (1.6pt);
  \node at ($(O)+(-0.24,-0.15)$) {$O$};

  % 顶点标注
  \node[below left] at (A) {$A$};
  \node[below right] at (B) {$B$};
  \node[right] at (C) {$C$};
  \node[left=1pt, yshift=2pt] at (D) {$D$};
  \node[left] at (A1) {$A_1$};
  \node[above=1pt] at (B1) {$B_1$};
  \node[right] at (C1) {$C_1$};
  \node[above left] at (D1) {$D_1$};
\end{tikzpicture}
\end{center}
\nopagebreak[4]
\end{question}
\begin{question}
\gkti{9}已知 $f(x)=\ln(1+x)-\ln(1-x)$，$x\in(-1,1)$。现有下列命题：

① $f(-x)=-f(x)$；② $f\left(\dfrac{2x}{x^2+1}\right)=2f(x)$；③ $|f(x)|\geqslant 2|x|$。

其中所有正确命题的序号是\mbox{（　　）}

（A）①②③；（B）②③；（C）①③；（D）①②。

\end{question}
\begin{question}
\gkti{10}已知 $F$ 是抛物线 $y^2=x$ 的焦点，点 $A$，$B$ 在该抛物线上且位于 $x$ 轴的两侧，$\overrightarrow{OA}\cdot\overrightarrow{OB}=2$（其中 $O$ 为坐标原点），则 $\triangle ABO$ 与 $\triangle AFO$ 面积之和的最小值是\mbox{（　　）}

（A）$2$；（B）$3$；（C）$\dfrac{17\sqrt2}8$；（D）$\sqrt{10}$。

\end{question}
\gkpart{二、填空题}

\begin{question}
\gkti{11}复数 $\dfrac{2-2\mathrm i}{1+\mathrm i}=$\gkblank{}。

\end{question}
\begin{question}
\gkti{12}设 $f(x)$ 是定义在 $\mathbf R$ 上的周期为 $2$ 的函数，当 $x\in[-1,1)$ 时，
\nopagebreak[4]
\[f(x)=\begin{cases}-4x^2+2,&-1\leqslant x<0,\\x,&0\leqslant x<1,\end{cases}\]
则 $f\left(\dfrac32\right)=$\gkblank{}。

\Needspace{0.4\textheight}
\end{question}
\begin{question}
\gkti{13}如图，从气球 $A$ 上测得正前方的河流的两岸 $B$，$C$ 的俯角分别为 $67^\circ$，$30^\circ$，此时气球的高是 $46\ \mathrm m$，则河流的宽度 $BC$ 约等于\gkblank{}$\mathrm m$。

（用四舍五入法将结果精确到个位。参考数据：$\sin67^\circ\approx 0.92$，$\cos67^\circ\approx 0.39$，$\sin37^\circ\approx 0.60$，$\cos37^\circ\approx 0.80$，$\sqrt3\approx 1.73$）
\nopagebreak[4]
\begin{center}% ============================================================
% f2014_sichuan_science_q13_river_width.tex
% 2014 四川卷（理科）第 13 题：气球上测河流宽度
% 来源 PDF：普通高考/2014/2014四川理.pdf
% 原页：第 1 页（A3 横排）中栏下部，印面 919，第 13 题题干下方
% 说明：A 为气球位置，其正下方为河岸上的垂足；B、C 为近岸与对岸
%       观测点，A 到 B、C 的俯角分别为 67°、30°；气球高 46 m；
%       图中给出过 A 的水平虚线、铅垂高度线与垂足处的直角标记。
% ============================================================
\begin{tikzpicture}[>=Stealth, line width=0.7pt, font=\small,
  x=1cm, y=1cm, line cap=round, line join=round]

  \coordinate (A) at (0,4);
  \coordinate (D) at (0,0);
  \coordinate (B) at (1.70,0);
  \coordinate (C) at (6.93,0);

  % 地面与过 A 的水平线
  \draw[dashed] (-0.80,0) -- (7.60,0);
  \draw[dashed] (-0.20,4) -- (2.60,4);

  % 铅垂线与视线
  \draw[dashed] (A) -- (D);
  \draw[dashed] (A) -- (B);
  \draw[dashed] (A) -- (C);

  % 气球
  \draw (0,4.46) circle (0.22);
  \draw (-0.08,4.27) -- (0.08,4.27) -- (0.10,4.00) -- (-0.10,4.00) -- cycle;
  \node[font=\footnotesize] at (0,4.46) {$A$};

  % 铅垂高度与标注
  \draw (-0.85,4) -- (-0.25,4);
  \draw (-0.85,0) -- (-0.25,0);
  \draw[<->] (-0.55,4) -- (-0.55,0);
  \node[left] at (-0.70,2) {$46\mathrm{m}$};

  % 垂足处的直角标记
  \draw (0,0.16) -- (0.16,0.16) -- (0.16,0);

  % 俯角：30° 到远岸 C，67° 到近岸 B
  \draw ($(A)+(0.90,0)$) arc (0:-30:0.90);
  \node at (1.22,3.74) {$30^\circ$};
  \draw ($(A)+(1.70,0)$) arc (0:-67:1.70);
  \node at (2.05,3.68) {$67^\circ$};

  % 点的标注
  \node[below] at (B) {$B$};
  \node[below] at (C) {$C$};
\end{tikzpicture}
\end{center}
\nopagebreak[4]
\end{question}
\begin{question}
\gkti{14}设 $m\in\mathbf R$，过定点 $A$ 的动直线 $x+my=0$ 和过定点 $B$ 的动直线 $mx-y-m+3=0$ 交于点 $P(x,y)$，则 $|PA|\cdot|PB|$ 的最大值是\gkblank{}。

\end{question}
\begin{question}
\gkti{15}以 $A$ 表示值域为 $\mathbf R$ 的函数组成的集合，$B$ 表示具有如下性质的函数 $\varphi(x)$ 组成的集合：对于函数 $\varphi(x)$，存在一个正数 $M$，使得函数 $\varphi(x)$ 的值域包含于区间 $[-M,M]$。例如，当 $\varphi_1(x)=x^3$，$\varphi_2(x)=\sin x$ 时，$\varphi_1(x)\in A$，$\varphi_2(x)\in B$。现有如下命题：

① 设函数 $f(x)$ 的定义域为 $D$，则“$f(x)\in A$”的充要条件是“$\forall b\in\mathbf R$，$\exists a\in D$，$f(a)=b$”；

② 函数 $f(x)\in B$ 的充要条件是 $f(x)$ 有最大值和最小值；

③ 若函数 $f(x)$，$g(x)$ 的定义域相同，且 $f(x)\in A$，$g(x)\in B$，则 $f(x)+g(x)\notin B$；

④ 若函数 $f(x)=a\ln(x+2)+\dfrac{x}{x^2+1}$（$x>-2$，$a\in\mathbf R$）有最大值，则 $f(x)\in B$。

其中的真命题有\gkblank{}。（写出所有真命题的序号）

\end{question}
\gkpart{三、解答题}

\begin{question}
\gkti{16}已知函数 $f(x)=\sin\left(3x+\dfrac\pi4\right)$。

（1） 求 $f(x)$ 的单调递增区间；

（2） 若 $\alpha$ 是第二象限角，$f\left(\dfrac\alpha3\right)=\dfrac45\cos\left(\alpha+\dfrac\pi4\right)\cos 2\alpha$，求 $\cos\alpha-\sin\alpha$ 的值。

\end{question}
\begin{question}
\gkti{17}一款击鼓小游戏的规则如下：每盘游戏都需要击鼓三次，每次击鼓要么出现一次音乐，要么不出现音乐；每盘游戏击鼓三次后，出现一次音乐获得 $10$ 分，出现两次音乐获得 $20$ 分，出现三次音乐获得 $100$ 分，没有出现音乐则扣除 $200$ 分（即获得 $-200$ 分）。设每次击鼓出现音乐的概率为 $\dfrac12$，且各次击鼓出现音乐相互独立。

（1） 设每盘游戏获得的分数为 $X$，求 $X$ 的分布列；

（2） 玩三盘游戏，至少有一盘出现音乐的概率是多少？

（3） 玩过这款游戏的许多人都发现，若干盘游戏后，与最初的分数相比，分数没有增加反而减少了。请运用概率统计的相关知识分析分数减少的原因。

\Needspace{0.4\textheight}
\end{question}
\begin{question}
\gkti{18}三棱锥 $A-BCD$ 及其侧视图、俯视图如图所示。设 $M$，$N$ 分别为线段 $AD$，$AB$ 的中点，$P$ 为线段 $BC$ 上的点，且 $MN\perp NP$。

（1） 证明：$P$ 为线段 $BC$ 的中点；

（2） 求二面角 $A-NP-M$ 的余弦值。
\nopagebreak[4]
\begin{center}% ============================================================
% f2014_sichuan_science_q18_tetrahedron.tex
% 2014 四川卷（理科）第 18 题：三棱锥及其侧视图、俯视图
% 来源 PDF：普通高考/2014/2014四川理.pdf
% 原页：第 2 页（A3 横排）左栏上部，印面 920，第 18 题题干下方
% 说明：左侧为三棱锥 A-BCD 的直观图（M、N 分别为 AD、AB 中点，
%       P 为 BC 中点）；中间为侧视图（等腰三角形，腰 2、底 2，高
%       线为虚线）；右侧为俯视图（直立边 2 被中点分为 1+1，两腰
%       均为 2，并连出中点与右顶点）。
% ============================================================
\begin{tikzpicture}[>=Stealth, line width=0.7pt, font=\small,
  x=1cm, y=1cm, line cap=round, line join=round]

  % ---- 左：三棱锥直观图 ----
  \coordinate (A) at (0,4.20);
  \coordinate (B) at (-0.58,1.09);
  \coordinate (C) at (3.07,1.63);
  \coordinate (D) at (0.74,2.37);
  \coordinate (N) at ($(A)!0.5!(B)$);
  \coordinate (M) at ($(A)!0.5!(D)$);
  \coordinate (P) at ($(B)!0.5!(C)$);

  \draw (A) -- (B) -- (C) -- cycle;
  \draw (A) -- (C);
  \draw[dashed] (A) -- (D);
  \draw[dashed] (D) -- (C);
  \draw[dashed] (D) -- (B);
  \draw[dashed] (N) -- (M);
  \draw (N) -- (P);

  \node[above] at (A) {$A$};
  \node[below left] at (B) {$B$};
  \node[right] at (C) {$C$};
  \node[above right] at (D) {$D$};
  \node[left] at (N) {$N$};
  \node[above right] at (M) {$M$};
  \node[below] at (P) {$P$};

  % ---- 中：侧视图 ----
  \begin{scope}[shift={(5.60,0.90)}]
    \coordinate (S1) at (0,0);
    \coordinate (S2) at (2,0);
    \coordinate (S3) at (1,1.732);
    \draw (S1) -- (S2) -- (S3) -- cycle;
    \draw[dashed] (S3) -- (1,0);
    \node[below] at (0.50,-0.14) {$1$};
    \node[below] at (1.50,-0.14) {$1$};
    \node[left] at (0.36,0.92) {$2$};
    \node[right] at (1.64,0.92) {$2$};
    \node[below] at (1,-0.55) {侧视图};
  \end{scope}

  % ---- 右：俯视图 ----
  \begin{scope}[shift={(8.80,1.70)}]
    \coordinate (T1) at (0,1);
    \coordinate (T2) at (0,-1);
    \coordinate (T3) at (1.732,0);
    \draw (T1) -- (T2) -- (T3) -- cycle;
    \draw (0,0) -- (T3);
    \node[left] at (0,0.52) {$1$};
    \node[left] at (0,-0.52) {$1$};
    \node[above left] at (0.86,0.72) {$2$};
    \node[below left] at (0.86,-0.72) {$2$};
    \node[below] at (0.87,-1.35) {俯视图};
  \end{scope}

\end{tikzpicture}
\end{center}
\nopagebreak[4]
\end{question}
\begin{question}
\gkti{19}设等差数列 $\{a_n\}$ 的公差为 $d$，点 $(a_n,b_n)$ 在函数 $f(x)=2^x$ 的图象上（$n\in\mathbf N^*$）。

（1） 若 $a_1=-2$，点 $(a_8,4b_7)$ 在函数 $f(x)$ 的图象上，求数列 $\{a_n\}$ 的前 $n$ 项和 $S_n$；

（2） 若 $a_1=1$，函数 $f(x)$ 的图象在点 $(a_2,b_2)$ 处的切线在 $x$ 轴上的截距为 $2-\dfrac1{\ln 2}$，求数列 $\left\{\dfrac{a_n}{b_n}\right\}$ 的前 $n$ 项和 $T_n$。

\end{question}
\begin{question}
\gkti{20}已知椭圆 $C:\dfrac{x^2}{a^2}+\dfrac{y^2}{b^2}=1$（$a>b>0$）的焦距为 $4$，其短轴的两个端点与长轴的一个端点构成正三角形。

（1） 求椭圆 $C$ 的标准方程；

（2） 设 $F$ 为椭圆 $C$ 的左焦点，$T$ 为直线 $x=-3$ 上任意一点，过 $F$ 作 $TF$ 的垂线交椭圆 $C$ 于点 $P$，$Q$。

① 证明：$OT$ 平分线段 $PQ$（其中 $O$ 为坐标原点）；

② 当 $\dfrac{|TF|}{|PQ|}$ 最小时，求点 $T$ 的坐标。

\end{question}
\begin{question}
\gkti{21}已知函数 $f(x)=\mathrm e^x-ax^2-bx-1$，其中 $a$，$b\in\mathbf R$，$\mathrm e=2.71828\cdots$ 为自然对数的底数。

（1） 设 $g(x)$ 是函数 $f(x)$ 的导函数，求函数 $g(x)$ 在区间 $[0,1]$ 上的最小值；

（2） 若 $f(1)=0$，函数 $f(x)$ 在区间 $(0,1)$ 内有零点，求 $a$ 的取值范围。

\end{question}